% Since LLMs are stochastic systems which output a probability over next word solution, in theory an LLM can generate any student reference solution by some random chance. Under this assumption, we extract the random noise used to create a given solution and seed an LLM with this noise under a counterfactual score scenario. By leveraging the stochastic nature of LLMs and controlling the decoding process, we produce personalized, actionable feedback that is closely aligned with each student's submission.

% this may often fall short, as resulting counterfactual generation systems may learn shortcuts which affect scoring models without providing meaningful differences. This differs significantly from the rubric-based setting where feedback not only is required but

% A central challenge in designing a counterfactual generation approach is balancing the tradeoff of task performance and original factual text similarity. Existing approaches address this tradeoff by training a masking model that identifies tokens that have maximum impact on performance


% In addressing this tradeoff, existing approaches are designed to favor task performance, optimizing domain-specific models that do not adapt to other domains, or are trained to condition on more generic control codes. (Tightly coupled the approach has minimality baked in as a joint objective or element of a search process, looking for minimal token flips which results in a new label.)



% question implies the existance of a alternative, counterfactual version of student work which would have recieved a different (often improved score). However, there are infinitely many alterantive student answers some of which 


% building off existing student knowledge and encourage critical thinking based off learned concepts "What would the graph look like if we change $a$ from 1 to 3?" 

4. Contributions
a. novel approach to counterfactual feedback generation, extensive evaluation of ordinal task, Our results show this approach provides a novel inference-only method with a simple tunable control signal with no training

What is the gap??

creating  a fundamental trade-off between task performance and 

Gaps
1. Underexplored for education tasks, there is work in feedback but little to nothing in counterfactual feedback
2. Existing approaches to CFG focus on changing broad categorizations: topic/correctness/toxicity. These categorizations are nominal, not 
3. SOTA identify-and-replace require training a task-specific model for token identification, joint-objective models require training a vector representation or entire model for the task (infeasible with growing sizes of generative LLMS)
4. Other methods use control task-specific generic control codes which do not port well to an ordinal setting

Thus we arrive with a inference time procedure which is enhanced by a general method for model alignment which would translate across tasks

3. Gaps in existing work (might be multiple paragraphs in this case for ed, counterfactual, gumbel)
Existing approaches to automated contrastive feedback follow two approaches: joint learning, mask-and-fill, or most recently prompting-based. (Maybe instead shift this into focus on changing a single nominal label is limited in scope)

By leveraging LLMs, it becomes possible to generate realistic and contextually relevant modifications to student work that align with rubric criteria. This approach not only reduces the manual burden on instructors but also provides students with personalized examples of how their submissions could be improved, thereby enhancing the feedback process and supporting deeper learning.

% Rubric-based grading is highly informative for learners as it quantifies the specific areas in which their work could be improved, guiding continued study or training.

% Rubric-based grading is convenient for instructors, allowing them to fairly evaluate work which is often diverse in a way that links back to a curriculum.

% Existing LLM-based approaches to automated counterfactual generation are broadly categorized into three categories groups: those which directly optimize an LLM for the counterfactual generation task, those which  prompt LLMs, and those which use an LLM in an identify-and-generate approach. 



Feedback can come in many forms: a binary grade to extensive writing




4. Contributions
a. novel approach to counterfactual feedback generation, extensive evaluation of ordinal task, Our results show this approach provides a novel inference-only method with a simple tunable control signal with no training

Task: Generating counterfactual student answer that achieve a desired rubric score but remain recognizably close to the student’s original answer

Enviornment: LLMS


\mh{Under construction, do not read!}
% (1) A description of the high-level domain/environment
% (2) A key problem that happens in this high-level domain
% (3) Current solutions to that key problem
% (4) A problem with the current solutions
% (5) A glimpse at how you solve the problem in (4)
% (6) An overview of your evaluation/experiments results

\mh{TODO Update terminology: contrastive feedback}

Counterfactual ideation and framing has a long history of studies across disciplines including philosophy, logic, history, and medicine. In construction a counterfactual scenario, we ask what would have happened if a prior condition was changed. The implications of the question can have a wide variety of consequence: from "Would my stomach have hurt if I didn't eat that last cookie?" to "Would the patient still be with us if I used a different medication?" Many applications of Natural Language Processing can also be seen under the lens of counterfactual conditioning. For example Machine Translation could broadly be seen as, "How would this thought be described if expressed in Mandarin instead of English". And computer science research explorations often ask experiments seek to answer "How would this model perform on the task if algorithm B was used in place of algorithm A?"

In the context of education, counterfactual thinking provides instructors strategies to build off of existing student knowledge or work to facilitate learning. For example, in teaching writing an instructor could comment: "If you had a stronger topic sentence this paragraph, it would have helped the reader understand your thesis." This feedback implies the existence of a alternative, counterfactual student answer which would demonstrate better mastery of a particular skill. Providing a demonstration of a counterfactual example is relatively common practice in teaching, but is often general, taking the form of answer keys or showcasing exemplary peer work for a given assignment. While general feedback is helpful, a counterfactual example which is more similar to the original student would be more helpful, as a student can more directly see how their work can be immediately improved and build off their existing ideas and interests. 

While in theory, an instructor could craft an extension of each individual students' work, in practice this is infeasible as it would require extensive time and effort.  Automating this process, or at least providing a framework for generating counterfactual student answers, could help bridge the gap between rubric-based feedback and actionable guidance tailored to individual student submissions. The generative and instruction-following capabilities of large language models (LLMs) offer a promising avenue for automating feedback and edits to student work (TODO need cites). However, the example work of LLMs often diverge from the reference work, even when explicitly instructed not to, introducing vocabulary or concepts which are often are unfamiliar to students. This can result in feedback that is less actionable or relatable, as students may struggle to connect the generated work to their own work. Ensuring that counterfactual answers remain close in style, vocabulary, and conceptual approach to the original student submission is therefore crucial for maximizing the educational value of these example work.

In our work, we explore an approach towards counterfactual example generation which leverages a computational phenomena to recover random noise information from a dataset of real student answers and reuses this noise as a control signal for counterfactual example generation process. 

Our research questions can be summarized as follows:

RQ1: Under what conditions does recovered Gumbel noise allow for a system that can produce counterfactual student answers that adheres to rubric guidelines while staying lexically similar to the original student answer?

RQ2: How does the recovered noise interact with the autoregressive sampling of an LLM and are there conditions under which the noise is more effective at regulating output?

RQ3: What information is carried by this sampling noise? Are there any inherent patterns of differences between noise values which result in high quality results and those that do not?

Our findings and contributions can be summarized as:

Our empirical investigations find that existing Hindsight recovery and reuse appoaches do not lead to a clear recovery signal, at least in the context of the task of counterfactual example generation.

However, we find that under specific conditions that recovered noise can be used to control output generation, however the signal completely dominates the control process.

Leveraging our findings, we introduce a modified sampling and recovery algorithm Beta-Hindsight sampling, which allows for a tunable tradeoff between reference similarity and task performance.

% Observed Phenomenon:

% - Recovered noise does not behave as a neutral sampling perturbation.
% - Without strong training or reference hints, generations are unstable or nonsensical
% - With strong training/reference hinting the recovered noise overwhelms model behavior producing verbatim replicas and counterfactuality collapses
% - With noise regularization a tuanble tradeoff emerges between task fidelity and reference proximity